\documentclass{article}

\usepackage{arxiv}

\usepackage[utf8]{inputenc}
\usepackage[T1]{fontenc}
\usepackage{hyperref}
\usepackage{url}
\usepackage{booktabs}
\usepackage{amsfonts}
\usepackage{amsmath}
\usepackage{nicefrac}
\usepackage{microtype}
\usepackage{cleveref}
\usepackage{graphicx}
\usepackage{natbib}
\usepackage{doi}
\usepackage{bm}

\graphicspath{{assets/figures/}{assets/others/}}

\title{Bayesian MCMC Inference Engine for Cloud Compute Billing Prediction: An AR($p$) Approach}

\author{
    Axl Roel Andaya, Christian Joseph Bunyi, Sean Higgins Lim \\
    De La Salle University \\
    \texttt{\{axl\_andaya,christian\_bunyi,sean\_higgins\_lim\}} \\ \texttt{@dlsu.edu.ph}
}

% Header overrides
\renewcommand{\shorttitle}{Bayesian MCMC Cloud Compute Billing Prediction}
\date{}

\begin{document}
\maketitle

% \begin{abstract}
% The abstract.
% \end{abstract}

% \keywords{Bayesian Inference \and MCMC \and Gibbs Sampler \and Autoregressive Model \and Cloud Billing}

\section{Introduction}
Explain the motivation behind Topic 10 (Cloud Compute Billing). Provide the justification for using a Bayesian approach, specifically highlighting the value of uncertainty quantification in cost forecasting compared to frequentist methods.

\section{Methodology}
This section describes the mathematical formulation of our AR($p$) framework and the subsequent logic of the Markov Chain Monte Carlo (MCMC) inference engine.

\subsection{Mathematical Framework}
Detail the AR($p$) process equation and the choice of semi-conjugate priors (Multivariate Normal for coefficients $\bm{\beta}$ and Inverse-Gamma for variance $\sigma^2$). Provide the derived full conditional distributions used in the sampling process.

\subsection{MCMC Engine Implementation}
Explain the iterative Gibbs sampling loop. Describe the structure of the R package core and how the sampling transitions between the conditional posteriors of $\bm{\beta}$ and $\sigma^2$.

\section{Results and Discussion}
We present the empirical results of our engine, starting with a ground-truth simulation and moving into detailed diagnostics and forecasting results.

\subsection{Ground Truth Simulation \& Verification}
Establish the known parameters used to generate the synthetic series. Demonstrate that the MCMC engine successfully recovers these "true" coefficients, verifying the mathematical correctness of the implementation.

\subsection{Diagnostics \& Convergence Analysis}
Provide trace plots and density plots for the parameters to show mixing and stationarity. Report the Effective Sample Size (ESS) to ensure the independence of iterations.

\subsection{Forecasting \& Uncertainty Quantification}
Generate a 30-day forecast. Analyze the posterior predictive distributions and the 95\% credible intervals to quantify the range of likely billing outcomes.

% \section{Conclusion}
% Summary of the project's contributions and potential areas for future refinement.

\bibliographystyle{unsrtnat}
\bibliography{references}

\end{document}
